(1) 光学腔的往返ABCD矩阵稳定时其本征值模为1，普通非退化稳定区满足 $\boxed{0<g_1g_2<1}$。严格共焦情形 $g_1g_2=0$ 为特殊边界，可有有限Rayleigh长度的Gaussian自再现模，不应简单判作无法共振。

(2) 纵横模式频率为 $\boxed{\nu_{qmn}=\frac c{2L}[q+(m+n+1)\arccos\sqrt{g_1g_2}/\pi]}$，其中横向阶数之和 $m+n$ 通过Gouy相位导致横模频移。采用理想近轴、无镜面色散约定，轴向整数 $q$ 取适当纵向分支。

(3) 共焦时 $\arccos0=\pi/2$，故 $\nu_{qmn}=c[q+(m+n+1)/2]/(2L)$。束腰在镜间中点，镜面波前曲率 $R_m=L/2+z_R^2/(L/2)$ 等于 $L$，因而 $\boxed{z_R=L/2}$、$\boxed{w_0^2=\lambda L/(2\pi)}$。增加两个横向模阶并将 $q$ 减1，频率保持不变。

(4) 镜距发生偏差时Gouy相位系数不再严格等于 $1/2$，原有纵横模简并通常分裂。高阶模式的光斑更大，有限凹镜孔径会使其衍射损耗增大；镜面形状误差也造成额外频移；若腔内加入色散介质，纵模光程与Gouy相位均须按折射率重新计算，因此本题公式仅适用于理想近轴、低损耗的空气谐振腔。
